\chapter{Curved Multicomponent FLRW Background and Distances}
\PhysicsContext{Background geometry and perturbations}{The first part fixes the common FLRW and perturbation conventions used later. It is broader than the dark-sector model and supplies the transport language in which model-specific effects are measured.}



The homogeneous model is extended to spatial curvature, a cosmological constant, multiple fluids, and standard distance observables.  The equations define a controlled background solver rather than a microscopic derivation of every component.

\allowdisplaybreaks

\section{Curved lapse-resolved minisuperspace action}
Let
\begin{equation}
\dd s^2=-N(t)^2\dd t^2+a(t)^2\gamma_{ij}^{(K)}\dd x^i\dd x^j,
\qquad
{}^{(3)}R[\gamma^{(K)}]=6K.
\end{equation}
For unit comoving volume,
\begin{equation}
S_{\mathrm{FLRW}}
=\int\left[
-\frac{3M_{\mathrm{Pl}}^2a\dot a^2}{N}
+3M_{\mathrm{Pl}}^2KaN
-M_{\mathrm{Pl}}^2\Lambda a^3N
-NH_{\mathrm P}(a)
\right]\dd t.
\end{equation}
Varying \(N\) gives
\begin{align}
0
&=
\frac{3M_{\mathrm{Pl}}^2a\dot a^2}{N^2}
+3M_{\mathrm{Pl}}^2Ka
-M_{\mathrm{Pl}}^2\Lambda a^3
-H_{\mathrm P}(a).
\end{align}
After division by \(3M_{\mathrm{Pl}}^2a^3\) and use of
\(H=\dot a/(Na)\), this becomes the Friedmann equation.
With
\begin{equation}
\rho(a):=\frac{H_{\mathrm P}(a)}{a^3},
\end{equation}
the lapse equation gives
\begin{equation}
\boxed{
H^2
=\frac{\rho}{3M_{\mathrm{Pl}}^2}
+\frac{\Lambda}{3}
-\frac{K}{a^2}.
}
\end{equation}
For total pressure \(p\),
\begin{align}
\dot H
&=-\frac{\rho+p}{2M_{\mathrm{Pl}}^2}+\frac{K}{a^2},\\
\frac{\ddot a}{a}
&=-\frac{\rho+3p}{6M_{\mathrm{Pl}}^2}+\frac{\Lambda}{3},\\
\dot\rho+3H(\rho+p)&=0.
\end{align}
In the gauge \(N=1\), differentiate the Friedmann equation:
\begin{align}
2H\dot H
&=
\frac{\dot\rho}{3M_{\mathrm{Pl}}^2}
+\frac{2KH}{a^2}.
\end{align}
Using \(\dot\rho=-3H(\rho+p)\) and dividing by \(2H\) gives the equation for \(\dot H\).  Then
\begin{align}
\frac{\ddot a}{a}
&=\dot H+H^2\\
&=
-\frac{\rho+p}{2M_{\mathrm{Pl}}^2}
+\frac{K}{a^2}
+\frac{\rho}{3M_{\mathrm{Pl}}^2}
+\frac{\Lambda}{3}
-\frac{K}{a^2},
\end{align}
which simplifies to the acceleration equation.
The symbolic identity
\begin{equation}
\frac{a}{2}\frac{\dd}{\dd a}
\left[
\frac{8\pi G}{3}\rho+\frac{\Lambda}{3}-\frac{K}{a^2}
\right]
=-4\pi G(\rho+p)+\frac{K}{a^2}
\end{equation}
is exact when
\begin{equation}
a\rho_{,a}=-3(\rho+p).
\end{equation}

\section{Constant-equation-of-state component library}
For noninteracting components
\begin{equation}
p_i=w_i\rho_i,
\end{equation}
\begin{align}
\dot\rho_i+3H(1+w_i)\rho_i&=0,\\
\rho_i(a)&=\rho_{i0}a^{-3(1+w_i)}.
\end{align}
Since \(\dot\rho_i=aH\,\dd\rho_i/\dd a\),
\begin{align}
\frac{\dd\rho_i}{\rho_i}
&=-3(1+w_i)\frac{\dd a}{a}.
\end{align}
Integration from \(a=1\) to \(a\) gives the stated power law.
With \(a_0=1\),
\begin{align}
\Omega_{i0}&:=\frac{\rho_{i0}}{3M_{\mathrm{Pl}}^2H_0^2},\\
\Omega_{\Lambda0}&:=\frac{\Lambda}{3H_0^2},\\
\Omega_{K0}&:=-\frac{K}{H_0^2},\\
1&=\sum_i\Omega_{i0}+\Omega_{K0}+\Omega_{\Lambda0},\\
E(a)^2
:=\frac{H(a)^2}{H_0^2}
&=\sum_i\Omega_{i0}a^{-3(1+w_i)}
+\Omega_{K0}a^{-2}+\Omega_{\Lambda0}.
\end{align}
Hence
\begin{equation}
E(z)^2
=\sum_i\Omega_{i0}(1+z)^{3(1+w_i)}
+\Omega_{K0}(1+z)^2+\Omega_{\Lambda0}.
\end{equation}
Special cases:
\begin{align}
\rho_r&\propto a^{-4},\\
\rho_m&\propto a^{-3},\\
\rho_K&\propto a^{-2},\\
\rho_\Lambda&=\mathrm{constant}.
\end{align}

\section{Conservative interacting fluids}
Let
\begin{equation}
\dot\rho_i+3H(\rho_i+p_i)=Q_i,
\qquad
\sum_iQ_i=0.
\end{equation}
Then
\begin{equation}
\boxed{
\dot\rho+3H(\rho+p)=0,
\qquad
\rho:=\sum_i\rho_i,
\quad
p:=\sum_i p_i.
}
\end{equation}
Summing the component equations gives
\begin{align}
\dot\rho+3H(\rho+p)
&=\sum_iQ_i
=0.
\end{align}
Energy transfer therefore changes the partition among components but not the total homogeneous continuity equation.
For the benchmark transfer
\begin{align}
Q_m&=-\Gamma\rho_m,\\
Q_r&=+\Gamma\rho_m,
\end{align}
\begin{align}
\dot\rho_m&=-3H\rho_m-\Gamma\rho_m,\\
\dot\rho_r&=-4H\rho_r+\Gamma\rho_m,\\
\dot\rho_m+\dot\rho_r
+3H\left(\rho_m+\rho_r+\frac13\rho_r\right)&=0.
\end{align}

\section{Background observables}
Define
\begin{align}
\bar\chi(z)&:=\int_0^z\frac{\dd z'}{E(z')},\\
D_H&:=\frac{c}{H_0}.
\end{align}
For a radial null ray,
\begin{equation}
0=-c^2\dd t^2+a(t)^2\dd\chi_{\rm com}^2.
\end{equation}
Using \(1+z=a^{-1}\) and
\(\dd t=-\dd z/[(1+z)H(z)]\) gives
\begin{equation}
\dd\chi_{\rm com}
=\frac{c\,\dd z}{H(z)}
=D_H\frac{\dd z}{E(z)}.
\end{equation}
Thus \(D_H\bar\chi(z)\) is the radial comoving distance before the curvature map is applied.
The transverse comoving distance is
\begin{equation}
D_M(z)=D_H\,\mathcal S_{\Omega_K}\!\left(\bar\chi(z)\right),
\end{equation}
where
\begin{equation}
\mathcal S_{\Omega_K}(x)
:=
\begin{cases}
\Omega_K^{-1/2}\sinh(\sqrt{\Omega_K}\,x),&\Omega_K>0,\\
x,&\Omega_K=0,\\
|\Omega_K|^{-1/2}\sin(\sqrt{|\Omega_K|}\,x),&\Omega_K<0.
\end{cases}
\end{equation}
Then
\begin{align}
D_A(z)&=\frac{D_M(z)}{1+z},\\
D_L(z)&=(1+z)D_M(z),\\
&\boxed{D_L(z)=(1+z)^2D_A(z)},\\
H_0t_{\mathrm L}(z)
&=\int_0^z\frac{\dd z'}{(1+z')E(z')},\\
H_0t_0
&=\int_0^1\frac{\dd a}{aE(a)}.
\end{align}
Exact benchmarks follow by direct integration.  For matter domination, \(E(a)=a^{-3/2}\) and \(E(z)=(1+z)^{3/2}\), so
\begin{align}
H_0t_0
&=\int_0^1a^{1/2}\dd a=\frac23,\\
\bar\chi(z)
&=\int_0^z(1+z')^{-3/2}\dd z'\\
&=2\left[1-(1+z)^{-1/2}\right].
\end{align}
For radiation domination, \(E(a)=a^{-2}\) and \(E(z)=(1+z)^2\), giving
\begin{align}
H_0t_0
&=\int_0^1a\,\dd a=\frac12,\\
\bar\chi(z)
&=\int_0^z\frac{\dd z'}{(1+z')^2}
=\frac{z}{1+z}.
\end{align}
For pure curvature, \(E(a)=a^{-1}\) and \(E(z)=1+z\), so
\begin{align}
H_0t_0
&=\int_0^1\dd a=1,\\
\bar\chi(z)
&=\int_0^z\frac{\dd z'}{1+z'}
=\log(1+z),\\
\frac{D_M}{D_H}&=\sinh\bar\chi
\qquad (\Omega_K=1).
\end{align}

\section{Projector background solver}
With normalized time \(\tau=H_0t\),
\begin{equation}
\frac{\dd a}{\dd\tau}=aE(a).
\end{equation}
The explicit midpoint update is
\begin{align}
k_1&=a_nE(a_n),\\
k_2&=\left(a_n+\frac{\Delta\tau}{2}k_1\right)
E\left(a_n+\frac{\Delta\tau}{2}k_1\right),\\
a_{n+1}&=a_n+\Delta\tau k_2.
\end{align}
For the curved radiation--matter--\(\Lambda\) benchmark,
\begin{equation}
p_a=2.024278.
\end{equation}
For conservative matter-to-radiation transfer,
\begin{equation}
p_{\mathrm{transfer}}=2.054690.
\end{equation}
The recorded residuals are
\begin{align}
\varepsilon_{\mathrm{distance\ benchmarks}}
&=8.881784\times10^{-16},\\
\varepsilon_{\mathrm{age\ benchmarks}}
&=0,\\
\varepsilon_{D_L-(1+z)^2D_A}
&=0,\\
\varepsilon_{\mathrm{total\ interacting\ continuity}}
&=4.440892\times10^{-16}.
\end{align}

\section{Open extensions}
\begin{nextsteps}
\item Derive a microscopic action that produces the exchange terms \(Q_i\) rather than inserting them phenomenologically.
\item Prove positivity, regularity, and global-existence conditions for the interacting background system.
\item Extend the projector background equations to anisotropic Bianchi geometries.
\item Determine whether nonsingular bouncing solutions exist and establish their dynamical stability when they do.
\item Calibrate the resulting background model against cosmological distance and expansion data.
\end{nextsteps}
